\documentclass[10pt]{article}

% ---------- Document Identity (update these only) ----------
\newcommand{\DocStage}{}
\newcommand{\DocTitle}{Your Road to Tier One SOF}
\newcommand{\ProgramName}{182-Week Civilian-to-Selection Readiness Pipeline}
\newcommand{\ProgramSubtitle}{}

% ---------- Page ----------
\usepackage[legalpaper,left=0.5in,right=0.5in,top=0.6in,bottom=0.45in]{geometry}

% ---------- Typography ----------
\usepackage[T1]{fontenc}
\usepackage[utf8]{inputenc}
\usepackage{libertinus}
\usepackage{microtype}
\usepackage{parskip}
\usepackage{ragged2e}
\usepackage{textcomp}
\AtBeginDocument{\color{black}}
\AtBeginDocument{\pagecolor{white}}
\AtBeginDocument{\justifying}

% ---------- Landscape pages ----------
\usepackage{pdflscape}

% ---------- Color ----------
\usepackage[table]{xcolor}
\definecolor{StageBlue}{HTML}{1F4E79}
\definecolor{StageBlueDark}{HTML}{163A5A}
\definecolor{LightGray}{HTML}{F2F4F7}
\definecolor{MidGray}{HTML}{D9DEE5}
\definecolor{RuleGray}{HTML}{C7CED8}
\definecolor{HeaderInk}{HTML}{000000}
\definecolor{Ink}{HTML}{000000}

% ---------- Utility ----------
\usepackage{enumitem}
\sloppy
\setlength{\emergencystretch}{2em}

% ---------- Boxes / UI ----------
\usepackage[most]{tcolorbox}

\tcbset{
  charterTitle/.style={
    enhanced,
    colback=StageBlue!4,
    colframe=StageBlue,
    boxrule=1.25pt,
    arc=3pt,
    left=16pt,right=16pt,top=14pt,bottom=14pt,
    borderline north={3.2pt}{0pt}{StageBlue},
    borderline west={4pt}{0pt}{StageBlue},
    borderline south={1.25pt}{0pt}{StageBlue},
    drop shadow={black!25!white},
  },
  charterRule/.style={
    enhanced,
    colback=LightGray,
    colframe=RuleGray,
    boxrule=0.85pt,
    arc=3pt,
    left=12pt,right=12pt,top=10pt,bottom=10pt,
    borderline west={3pt}{0pt}{StageBlue},
  },
  charterBadge/.style={
    enhanced,
    colback=StageBlue,
    coltext=white,
    colframe=StageBlueDark,
    boxrule=0.6pt,
    arc=2pt,
    left=6pt,right=6pt,top=3pt,bottom=3pt,
  }
}
\newtcbox{\Badge}{nobeforeafter,charterBadge}

% ---------- Lists ----------
\setlist[itemize]{leftmargin=1.5em, itemsep=0.25em, topsep=0.25em}
\setlist[enumerate]{leftmargin=1.8em, itemsep=0.25em, topsep=0.25em}

% ---------- TOC ----------
\usepackage{tocloft}
\setcounter{tocdepth}{3}
\setlength{\cftbeforesecskip}{3pt}
\setlength{\cftbeforesubsecskip}{2pt}
\setlength{\cftbeforesubsubsecskip}{0pt}
\renewcommand{\cftsecfont}{\bfseries}
\renewcommand{\cftsecpagefont}{\bfseries}
\renewcommand{\cftsubsecfont}{\normalfont}
\renewcommand{\cftsubsecpagefont}{\normalfont}
\renewcommand{\cftsubsubsecfont}{\normalfont}
\renewcommand{\cftsubsubsecpagefont}{\normalfont}
\setlength{\cftsecindent}{0pt}
\setlength{\cftsubsecindent}{1.25em}
\setlength{\cftsubsubsecindent}{2.8em}
\setlength{\cftsecnumwidth}{2.1em}
\setlength{\cftsubsecnumwidth}{2.8em}
\setlength{\cftsubsubsecnumwidth}{3.6em}

% Remove the built-in TOC heading so only the custom heading is shown
\makeatletter
\renewcommand{\tableofcontents}{\@starttoc{toc}}
\makeatother

% ---------- Header/Footer: page number only ----------
\usepackage{fancyhdr}
\pagestyle{fancy}
\fancyhf{}
\renewcommand{\headrulewidth}{0pt}
\renewcommand{\footrulewidth}{0pt}
\fancyfoot[C]{\thepage}

% ---------- Section formatting ----------
\usepackage{titlesec}

\titleformat{\section}
  {\Large\bfseries\color{Ink}}
  {\thesection}{0.6em}{}
  [\vspace{0.25em}\textcolor{StageBlue}{\rule{\linewidth}{0.8pt}}]

\titleformat{\subsection}
  {\large\bfseries\color{Ink}}
  {\thesubsection}{0.6em}{}

\titleformat{\subsubsection}
  {\normalsize\bfseries\color{Ink}}
  {\thesubsubsection}{0.6em}{}

\titlespacing*{\section}{0pt}{0.95em}{0.35em}
\titlespacing*{\subsection}{0pt}{0.65em}{0.25em}
\titlespacing*{\subsubsection}{0pt}{0.55em}{0.2em}

% ---------- Table engine ----------
\usepackage{tabularray}
\UseTblrLibrary{booktabs}

% ---------- tabularray: GLOBAL table treatment ----------
\SetTblrInner{
  cells = {fg=black, bg=white, valign=t},
}

% ---------- Header helpers ----------
\newcommand{\Hdr}[1]{\shortstack[c]{\strut #1\strut}}

% ---------- Lines ----------
\newcommand{\TitleLine}{\textcolor{StageBlue}{\rule{0.98\linewidth}{0.95pt}}}
\newcommand{\SubTitleLine}{\textcolor{RuleGray}{\rule{0.98\linewidth}{0.65pt}}}

% ---------- Continued on next page ----------
\newcommand{\ContinuedOnNextPageInline}{%
  \par
  \vfill
  \noindent\textcolor{RuleGray}{\rule{\linewidth}{1pt}}\vspace{-2pt}
  \noindent\hfill\textit{Continued on next page}\par\vspace{-10pt}
  \noindent\textcolor{RuleGray}{\rule{\linewidth}{1pt}}
  \clearpage
}

% ---------- PDF hyperlinks ----------
\usepackage[hidelinks]{hyperref}
\hypersetup{
  colorlinks=false,
  pdfborder={0 0 0},
  linktoc=all,
  pdftitle={\DocTitle},
  pdfauthor={\ProgramName}
}

% =========================================================
\begin{document}

\section*{Stage 2.B — Metrics and Test Inventory}
\phantomsection
\addcontentsline{toc}{section}{Stage 2.B — Metrics and Test Inventory}

\subsection*{2.B.1 Metric \textbf{ID} Standard and Naming Rules}
\phantomsection
\addcontentsline{toc}{subsection}{2.B.1 Metric \textbf{ID} Standard and Naming Rules}

\subsubsection*{Metric \textbf{ID} format and retirement rule (mandatory):}
\phantomsection

\begin{itemize}
\item Metric IDs \textbf{MUST} use the stage-scoped format \textbf{MET-2B}-\#\#\# (three digits, sequential).
\item IDs \textbf{MUST} be assigned once and \textbf{MUST} NOT be reused, even if a metric is later deprecated, merged, or retired.
\item If an \textbf{ID} is ever retired, the retirement \textbf{MUST} be explicitly recorded in the change log and the \textbf{ID} \textbf{MUST} remain permanently reserved (non-reusable).
\item One metric equals one row equals one stable \textbf{ID}. Bundling unrelated constructs into one \textbf{ID} is prohibited.
\end{itemize}

\subsubsection*{Domain tagging rule (mandatory):}
\phantomsection

\begin{itemize}
\item Every metric \textbf{MUST} have exactly one Primary Domain Tag.
\item Primary Domain Tags \textbf{MUST} be selected only from the canonical set defined in Stage 2.A: \textbf{RUN}, \textbf{RUCK}, \textbf{SWIM}, \textbf{CAL}, \textbf{STR}, \textbf{BODYCOMP}, \textbf{DURABILITY}, \textbf{RECOVERY}, \textbf{AER}.
\item Non-canonical tags, synonyms, or alternative labels (including “\textbf{STRENGTH}” as a tag label) \textbf{MUST} NOT be used in Stage 2.B.
\item \textbf{AER} is the canonical domain for low-impact aerobic and non-run conditioning constructs (bike, rower, elliptical, step tests, walk-based aerobic checks) and \textbf{MUST} NOT be treated as interchangeable evidence for \textbf{RUN} tolerance or \textbf{RUN} performance.
\end{itemize}

\subsubsection*{Measurement type definitions (mandatory enum):}
\phantomsection

\begin{itemize}
\item \textbf{TEST}: discrete, standardized performance outcome intended to be produced under protocol control.
\item \textbf{TREND}: longitudinal tracking used to evaluate change over time under consistent method/conditions posture.
\item \textbf{READINESS}: decision-support markers used to interpret safety, tolerance, and evidence admissibility posture.
\item \textbf{COMPLIANCE}: governance compliance markers used to compute admissible execution rates and enforce operating constraints.
\end{itemize}

\subsubsection*{Gate-Critical designation rule (locked defaults):}
\phantomsection

\begin{itemize}
\item Gate-Critical \textbf{MUST} be set to Yes for:
  \begin{itemize}
  \item all Stage 0.3 end-state target metrics, and
  \item the required compliance and risk-control metrics:
    \begin{itemize}
    \item Session Validity Tag,
    \item Weekly Admissible Session Compliance Percent,
    \item Cardio Minimum Standard Met,
    \item gait-altering pain flag.
    \end{itemize}
  \end{itemize}
\item All other metrics \textbf{MUST} default to Supporting (Gate-Critical = No) unless later elevated by a phase exit standard or gate definition under controlled change.
\end{itemize}

\subsubsection*{Do and Do Not (enforceable):}
\phantomsection

Do:

\begin{itemize}
\item Do keep metric names stable and descriptive; names \textbf{MUST} remain comparable across time.
\item Do keep units explicit and stable; any unit change requires controlled change with explicit comparability impact.
\item Do keep definitions operational and auditable; each metric \textbf{MUST} be interpretable by a third-party reviewer.
\item Do reference an associated protocol card C\# where obvious; use \textbf{TBD} 2.C only when protocol assignment is not yet determined.
\end{itemize}

Do Not:

\begin{itemize}
\item Do Not rename a metric to reset baselines, hide drift, or obscure poor comparability.
\item Do Not change units silently.
\item Do Not create duplicates (two different IDs for the same construct) to inflate progress narratives.
\item Do Not use non-canonical domain tags or “close enough” tag substitutions.
\end{itemize}

Every \textbf{MET-2B}-\#\#\# row \textbf{MUST} be interpretable without oral context. Each metric \textbf{MUST} be defined with enough structure that a third-party reviewer can answer, at minimum:

\begin{itemize}
\item What is being measured (construct definition).
\item How it is expressed (unit, formatting).
\item Under what conditions it is admissible for gate use (high-level validity requirements).
\item How it is collected (collection cadence high-level).
\item Whether it is gate-critical (Yes/No).
\item Which protocol card(s) govern(s) execution and admissibility (C\# reference or \textbf{TBD} 2.C in Draft — Non-Deployable posture only; use N/A when the metric is not executed as a protocolized test).
\item What it is used for (primary use statement).
\end{itemize}

\subsubsection*{Required interpretability posture:}
\phantomsection

\begin{itemize}
\item A metric name is not a definition. The definition is carried by the combination of Domain Tag + Measurement Type + Unit + Validity Requirements + Protocol Card reference(s).
\item If two metrics share the same construct, they must be merged under one \textbf{ID}; if they measure meaningfully different constructs, they must remain separate IDs.
\end{itemize}

\begin{itemize}
\item The Domain Tag column in Stage 2.B stores the Primary Domain Tag only.
\item Secondary domain tagging (when used) is not recorded in the Stage 2.B primary inventory table. If needed, it is recorded in:
  \begin{itemize}
  \item protocol card validity and disqualifier notes, and/or
  \item Stage 2.E logging conventions,
  \end{itemize}
  without creating dual ownership in the metrics inventory.
\end{itemize}

\subsubsection*{Unit locking (mandatory):}
\phantomsection

\begin{itemize}
\item A metric’s unit is part of its definition. Changing a unit is a comparability break and requires controlled change with explicit impact documentation.
\item Units must be recorded in a consistent display format across logs and reports.
\end{itemize}

\subsubsection*{Formatting rules (mandatory):}
\phantomsection

\begin{itemize}
\item Time-based metrics expressed as seconds (mm:ss):
  \begin{itemize}
  \item Primary stored value is total seconds.
  \item Display format may include mm:ss, but conversion must be consistent and traceable.
  \end{itemize}
\item Bodyweight expressed as lb (kg):
  \begin{itemize}
  \item Primary stored value may be lb with optional kg in parentheses, or vice versa, but the chosen storage posture must remain consistent across the program.
  \end{itemize}
\item Circumference expressed as in (cm):
  \begin{itemize}
  \item Sites must be stable once adopted.
  \item Tape method and posture must remain consistent within trend windows.
  \end{itemize}
\end{itemize}

\subsubsection*{Directionality (mandatory definition posture):}
\phantomsection

\begin{itemize}
\item Time trials: lower is better.
\item Rep counts: higher is better.
\item Holds (seconds): higher is better.
\item Flags (Y/N): Y indicates presence; for gate logic, the meaning of “Y” must be specified (disqualifier, warning, or informational).
\item 0–10 scales: higher indicates greater severity unless explicitly stated otherwise; interpretation remains conservative.
\end{itemize}

\subsubsection*{The Validity Requirements High-Level field \textbf{MUST} include, at a minimum:}
\phantomsection

\begin{itemize}
\item Standardization posture (what must remain stable for comparability).
\item Minimum record requirements (what must be logged to be admissible).
\item Explicit invalidity triggers (what causes the evidence to be treated as non-admissible for gate use).
\end{itemize}

\subsubsection*{High-level validity requirements must remain protocol-agnostic in Stage 2.B:}
\phantomsection

\begin{itemize}
\item Stage 2.B does not write stepwise procedures; it defines what conditions must be recorded and what breaks admissibility.
\item Detailed procedures and exact validity condition field forms are owned by Stage 2.C and Stage 2.E.
\end{itemize}

\begin{itemize}
\item Gate-Critical = Yes means:
  \begin{itemize}
  \item the metric is required for decision-grade gate interpretation for at least one gate posture, and/or
  \item the metric is a governance disqualifier/eligibility control for gate decisions.
  \end{itemize}
\item Gate-Critical metrics cannot be “skipped without consequence.” If missing or invalid, default posture is conservative:
  \begin{itemize}
  \item gate-required tests: default \textbf{HOLD} until valid evidence exists,
  \item disqualifier flags: default \textbf{HOLD}/\textbf{REGRESS}/\textbf{MEDICAL} \textbf{HOLD} as specified upstream.
  \end{itemize}
\item Gate-Critical = No means:
  \begin{itemize}
  \item the metric is supporting evidence only and does not independently justify \textbf{ADVANCE}.
  \item Supporting metrics can still constrain gates when they indicate safety or validity compromise (for example, supporting durability markers).
  \end{itemize}
\end{itemize}

\subsubsection*{Retirement rules (mandatory):}
\phantomsection

\begin{itemize}
\item A retired metric remains in the inventory with its \textbf{ID} reserved.
\item Retirement record must include:
  \begin{itemize}
  \item Metric \textbf{ID},
  \item retirement date,
  \item reason (duplicate, merged, out-of-scope, superseded),
  \item replacement metric \textbf{ID} if merged/superseded,
  \item comparability posture for pre- and post-retirement evidence.
  \end{itemize}
\end{itemize}

\subsubsection*{Merge rules (mandatory):}
\phantomsection

\begin{itemize}
\item If two metrics are found to measure the same construct, one \textbf{ID} is retired and mapped to the surviving \textbf{ID}.
\item Historical evidence is not rewritten; it is mapped through the retirement record.
\end{itemize}

\subsubsection*{Deprecation posture (permitted):}
\phantomsection

\begin{itemize}
\item A metric may be retained but marked “Supporting only” by policy, without changing its \textbf{ID}, to reduce misuse.
\item Deprecation must not be used to hide poor outcomes; it must be tied to construct validity and crosswalk clarity.
\end{itemize}

\begin{itemize}
\item \textbf{TBD} 2.C is permitted only while Stage 2 remains Draft — Non-Deployable during protocol-card authorship.
\item Gate-Critical metrics \textbf{MUST} NOT remain \textbf{TBD} 2.C at the point Stage 2 is declared Deployable — Gate Passed.
\item \textbf{TBD} 2.C is not permission to execute tests without validity posture; it is a placeholder indicating protocol assignment is incomplete.
\item For metrics that are not executed as protocolized tests (for example, derived \textbf{COMPLIANCE} metrics computed from logs), Associated Protocol Card C\# \textbf{MUST} be recorded as N/A (not \textbf{TBD} 2.C).
\end{itemize}

\begin{itemize}
\item None.
\end{itemize}

\begin{itemize}
\item Conflict source: baseline materials include legacy metric \textbf{ID} formats (for example, \textbf{MET}-\#\#\# and domain-prefixed \textbf{MET}-\#\#\# labels) that do not conform to the required \textbf{MET-2B}-\#\#\# convention. Adopted: \textbf{MET-2B}-\#\#\# is authoritative for Stage 2.B. Overridden: legacy \textbf{ID} formats. Why: this prompt locks Stage 2.B identifier conventions as crosswalk keys.
\item Conflict source: baseline materials do not consistently include the required non-run conditioning domain construct. Adopted: \textbf{AER} is included as a canonical domain tag for non-run conditioning constructs. Overridden: any prior tendency to classify low-impact machine conditioning as \textbf{RUN}. Why: prompt explicitly requires \textbf{AER} to prevent evidence contamination and unsafe impact equivalence.
\item Conflict source: baseline materials sometimes use “\textbf{STRENGTH}” or other synonyms as a domain tag label. Adopted: \textbf{STR} is canonical for strength and external load capacity. Overridden: the tag \textbf{STRENGTH} and any other non-canonical synonyms. Why: Stage 2.A canonical tags define \textbf{STR}; Stage 2.B prohibits non-canonical domain labels.
\end{itemize}

\clearpage
\begin{landscape}

\subsection*{2.B.2 Master Metrics Inventory}
\phantomsection
\addcontentsline{toc}{subsection}{2.B.2 Master Metrics Inventory}

\begingroup
\scriptsize
\setlength{\tabcolsep}{2pt}
\renewcommand{\arraystretch}{1.12}

\begin{longtblr}{
  width=\linewidth,
  rowhead=1,
  colspec = {
    Q[l,wd=0.55in]
    X[l,1.55]
    Q[l,wd=0.55in]
    Q[l,wd=0.75in]
    Q[l,wd=0.65in]
    Q[l,wd=0.95in]
    Q[l,wd=0.75in]
    Q[l,wd=0.95in]
    X[l,2.05]
    X[l,1.35]
  },
  hlines,
  vlines,
  cells = {hsep=6pt, vsep=6pt, halign=l, valign=t},
  cell{1-Z}{1-10} = {fg=black},
  row{odd} = {bg=white},
  row{even} = {bg=LightGray},
  row{1} = {bg=StageBlue!15, fg=black, font=\bfseries, halign=l, valign=t},
}
\Hdr{Metric \textbf{ID}} &
\Hdr{Metric Name} &
\Hdr{Domain} &
\Hdr{Measurement} &
\Hdr{Unit} &
\Hdr{Collection Cadence} &
\Hdr{Gate-Critical} &
\Hdr{Associated Protocol} &
\Hdr{Validity Requirements} &
\Hdr{Primary Use} \\
\textbf{MET-2B-001} & Bodyweight — Weekly Average & \textbf{BODYCOMP} & \textbf{TREND} & lb (kg) & weekly & Yes & \textbf{C1} & Consistent scale conditions; method/conditions recorded; no retroactive rewriting; anomalies noted & End-state target evidence; body composition trend governance \\
\textbf{MET-2B-002} & Body Fat Estimate Percent — Method/Device Recorded & \textbf{BODYCOMP} & \textbf{TREND} & \% (estimate) & monthly & Yes & \textbf{C2} & Method/device and conditions recorded each time; method consistency required within trend windows; no medical claims & End-state target evidence; trend comparability control \\
\textbf{MET-2B-003} & Circumference — Neck & \textbf{BODYCOMP} & \textbf{TREND} & in (cm) & weekly & No & \textbf{C3} & Site definition fixed; consistent tape and posture; conditions recorded; method drift logged & Supporting body composition trend and audit trace \\
\textbf{MET-2B-004} & Circumference — Chest & \textbf{BODYCOMP} & \textbf{TREND} & in (cm) & weekly & No & \textbf{C3} & Site definition fixed; consistent tape and posture; conditions recorded; method drift logged & Supporting body composition trend and audit trace \\
\textbf{MET-2B-005} & Circumference — Bicep & \textbf{BODYCOMP} & \textbf{TREND} & in (cm) & weekly & No & \textbf{C3} & Site definition fixed; consistent tape and posture; conditions recorded; method drift logged & Supporting body composition trend and audit trace \\
\textbf{MET-2B-006} & Circumference — Forearm & \textbf{BODYCOMP} & \textbf{TREND} & in (cm) & weekly & No & \textbf{C3} & Site definition fixed; consistent tape and posture; conditions recorded; method drift logged & Supporting body composition trend and audit trace \\
\textbf{MET-2B-007} & Circumference — Waist & \textbf{BODYCOMP} & \textbf{TREND} & in (cm) & weekly & No & \textbf{C3} & Site definition fixed; consistent tape and posture; conditions recorded; method drift logged & Supporting body composition trend and audit trace \\
\textbf{MET-2B-008} & Circumference — Thigh & \textbf{BODYCOMP} & \textbf{TREND} & in (cm) & weekly & No & \textbf{C3} & Site definition fixed; consistent tape and posture; conditions recorded; method drift logged & Supporting body composition trend and audit trace \\
\textbf{MET-2B-009} & Circumference — Calf & \textbf{BODYCOMP} & \textbf{TREND} & in (cm) & weekly & No & \textbf{C3} & Site definition fixed; consistent tape and posture; conditions recorded; method drift logged & Supporting body composition trend and audit trace \\
\textbf{MET-2B-010} & Push-Ups — 2-Minute Timed, Strict & \textbf{CAL} & \textbf{TEST} & reps & test weeks & Yes & \textbf{C4} & Standard conditions; warm-up completion logged; strict standard applied; no stop-rule violations; context recorded & End-state target evidence; readiness tier classification input \\
\textbf{MET-2B-011} & Sit-Ups — 2-Minute Timed, Standard & \textbf{CAL} & \textbf{TEST} & reps & test weeks & Yes & \textbf{C5} & Standard conditions; warm-up completion logged; standard applied; no stop-rule violations; context recorded & End-state target evidence; readiness tier classification input \\
\textbf{MET-2B-012} & Pull-Ups — Strict Dead-Hang & \textbf{CAL} & \textbf{TEST} & reps & test weeks & Yes & \textbf{C6} & Standard conditions; warm-up completion logged; strict dead-hang standard; no stop-rule violations; context recorded & End-state target evidence; readiness tier classification input \\
\textbf{MET-2B-013} & Plank — Forearm & \textbf{CAL} & \textbf{TEST} & seconds & test weeks & Yes & \textbf{C7} & Standard conditions; warm-up logged; standard form rules; termination on form failure; context recorded & End-state target evidence; readiness tier classification input \\
\textbf{MET-2B-014} & 1.5-Mile Run Time Trial & \textbf{RUN} & \textbf{TEST} & seconds (mm:ss) & test weeks & Yes & \textbf{C8} & Measured route OR treadmill admissible when device identity and settings are logged; warm-up logged; no stop-rule violations; environment recorded & End-state target evidence; readiness tier classification input \\
\textbf{MET-2B-015} & 2-Mile Run Time Trial & \textbf{RUN} & \textbf{TEST} & seconds (mm:ss) & test weeks & Yes & \textbf{C9} & Measured route OR treadmill admissible when device identity and settings are logged; warm-up logged; no stop-rule violations; environment recorded & End-state target evidence; readiness tier classification input \\
\textbf{MET-2B-016} & 3-Mile Run Time Trial & \textbf{RUN} & \textbf{TEST} & seconds (mm:ss) & test weeks & Yes & \textbf{C10} & Measured route OR treadmill admissible when device identity and settings are logged; warm-up logged; no stop-rule violations; environment recorded & End-state target evidence; readiness tier classification input \\
\textbf{MET-2B-017} & 5-Mile Run Time Trial & \textbf{RUN} & \textbf{TEST} & seconds (mm:ss) & test weeks & Yes & \textbf{C11} & Measured route OR treadmill admissible when device identity and settings are logged; warm-up logged; no stop-rule violations; environment recorded & End-state target evidence; readiness tier classification input \\
\textbf{MET-2B-018} & 3-Mile Ruck Time Trial — 25 lb & \textbf{RUCK} & \textbf{TEST} & seconds (mm:ss) & test weeks & Yes & \textbf{C12} & Measured route OR treadmill admissible when device identity and settings are logged; environment recorded; load and footwear recorded; warm-up logged; no stop-rule violations & End-state target evidence; readiness tier classification input \\
\textbf{MET-2B-019} & 6-Mile Ruck Time Trial — 35 lb & \textbf{RUCK} & \textbf{TEST} & seconds (mm:ss) & test weeks & Yes & \textbf{C12} & Measured route OR treadmill admissible when device identity and settings are logged; environment recorded; load and footwear recorded; warm-up logged; no stop-rule violations & End-state target evidence; readiness tier classification input \\
\textbf{MET-2B-020} & 9-Mile Ruck Time Trial — 45 lb & \textbf{RUCK} & \textbf{TEST} & seconds (mm:ss) & test weeks & Yes & \textbf{C12} & Measured route OR treadmill admissible when device identity and settings are logged; environment recorded; load and footwear recorded; warm-up logged; no stop-rule violations & End-state target evidence; readiness tier classification input \\
\textbf{MET-2B-021} & 12-Mile Ruck Time Trial — 55 lb & \textbf{RUCK} & \textbf{TEST} & seconds (mm:ss) & test weeks & Yes & \textbf{C12} & Measured route OR treadmill admissible when device identity and settings are logged; environment recorded; load and footwear recorded; warm-up logged; no stop-rule violations & End-state target evidence; readiness tier classification input \\
\textbf{MET-2B-022} & Swim 500 yd Time — No Fins & \textbf{SWIM} & \textbf{TEST} & seconds (mm:ss) & test weeks & Yes & \textbf{C13} & Unit locked (yards); no fins confirmed; warm-up logged; pool verification logged; no stop-rule violations & End-state target evidence; readiness tier classification input \\
\textbf{MET-2B-023} & Swim 500 m Time — No Fins & \textbf{SWIM} & \textbf{TEST} & seconds (mm:ss) & test weeks & Yes & \textbf{C13} & Unit locked (meters); no fins confirmed; warm-up logged; pool verification logged; no stop-rule violations & End-state target evidence; readiness tier classification input \\
\textbf{MET-2B-024} & Surface Water Confidence Exposure Flag & \textbf{SWIM} & \textbf{COMPLIANCE} & Y/N & per session & No & \textbf{C13} & Governance marker only; surface-only posture unless governed; environment/supervision context recorded when applicable & Supporting aquatic exposure trace and governance oversight \\
\textbf{MET-2B-025} & Underwater Exposure Admissibility Flag & \textbf{SWIM} & \textbf{COMPLIANCE} & Y/N & per underwater test event & Yes & \textbf{C14} & \textbf{MUST} be supervised and governance-permitted; certified supervision requirement; unsupervised attempt prohibited; governance prerequisites recorded & End-state underwater governance admissibility control \\
\textbf{MET-2B-026} & Underwater 25 m — One-Breath Time & \textbf{SWIM} & \textbf{TEST} & seconds (mm:ss) & test weeks & Yes & \textbf{C14} & Governance-constrained: certified supervision; rescue-ready posture; abort discipline; pool verification; unsupervised attempt prohibited & End-state target evidence (admissible only under governance) \\
\textbf{MET-2B-027} & Hinge Strength Proxy — \textbf{3-RM} / \textbf{5-RM} Load & \textbf{STR} & \textbf{TEST} & lb (kg) & test weeks & No & \textbf{C17} & Standardized implement/setup; warm-up logged; measurement method recorded; no stop-rule violations; submax risk posture enforced per protocol & Supporting strength progression classification (non-end-state target) \\
\textbf{MET-2B-028} & Squat Strength Proxy — \textbf{3-RM} / \textbf{5-RM} Load & \textbf{STR} & \textbf{TEST} & lb (kg) & test weeks & No & \textbf{C17} & Standardized implement/setup; warm-up logged; measurement method recorded; no stop-rule violations; submax risk posture enforced per protocol & Supporting strength progression classification (non-end-state target) \\
\textbf{MET-2B-029} & Upper Push Strength Proxy — \textbf{3-RM} / \textbf{5-RM} Load & \textbf{STR} & \textbf{TEST} & lb (kg) & test weeks & No & \textbf{C17} & Standardized implement/setup; warm-up logged; measurement method recorded; no stop-rule violations; submax risk posture enforced per protocol & Supporting strength progression classification (non-end-state target) \\
\textbf{MET-2B-030} & Upper Pull Strength Proxy — \textbf{3-RM} / \textbf{5-RM} Load & \textbf{STR} & \textbf{TEST} & lb (kg) & test weeks & No & \textbf{C17} & Standardized implement/setup; warm-up logged; measurement method recorded; no stop-rule violations; submax risk posture enforced per protocol & Supporting strength progression classification (non-end-state target) \\
\textbf{MET-2B-031} & Carry Capacity — Load Outcome & \textbf{STR} & \textbf{TEST} & lb (kg) & test weeks & No & \textbf{C17} & Carry setup standardized per protocol; warm-up logged; safe execution; no stop-rule violations; context recorded & Supporting load-handling capability classification (non-ruck construct) \\
\textbf{MET-2B-032} & Rower Time Trial & \textbf{AER} & \textbf{TEST} & seconds (mm:ss) & test weeks & No & \textbf{C18} & Machine identity/settings recorded; warm-up logged; pauses recorded; no stop-rule violations; context recorded & Supporting low-impact conditioning performance construct \\
\textbf{MET-2B-033} & Bike Time Trial & \textbf{AER} & \textbf{TEST} & seconds (mm:ss) & test weeks & No & \textbf{C18} & Machine identity/settings recorded; warm-up logged; pauses recorded; no stop-rule violations; context recorded & Supporting low-impact conditioning performance construct \\
\textbf{MET-2B-034} & Elliptical Time Trial & \textbf{AER} & \textbf{TEST} & seconds (mm:ss) & test weeks & No & \textbf{C18} & Machine identity/settings recorded; warm-up logged; pauses recorded; no stop-rule violations; context recorded & Supporting low-impact conditioning performance construct \\
\textbf{MET-2B-035} & Pain Severity — 0–10 + Location Recorded & \textbf{DURABILITY} & \textbf{READINESS} & 0–10 & daily & No & \textbf{C16} & Same-day capture; location required; trend interpreted conservatively; non-diagnostic posture & Supporting risk and tolerance interpretation \\
\textbf{MET-2B-036} & Gait-Altering Pain Flag & \textbf{DURABILITY} & \textbf{COMPLIANCE} & Y/N & per session & Yes & \textbf{C16} & If Yes, treat as disqualifier signal; documentation required; stop-rule posture honored; no concealment & Gate disqualifier signal and regression constraint \\
\textbf{MET-2B-037} & Foot Hotspot or Blister Flag + Location & \textbf{DURABILITY} & \textbf{COMPLIANCE} & Y/N & per session & No & \textbf{C16} & Same-day capture; location required; documentation required; used to constrain impact/load escalation & Supporting foot-interface risk governance \\
\textbf{MET-2B-038} & Sleep Hours & \textbf{RECOVERY} & \textbf{READINESS} & hours & daily & No & \textbf{C15} & Same-day capture; self-report permitted; trend interpreted conservatively for validity and risk & Supporting readiness interpretation for validity and safety \\
\textbf{MET-2B-039} & Sleep Quality & \textbf{RECOVERY} & \textbf{READINESS} & 0–10 & daily & No & \textbf{C15} & Same-day capture; scale anchored consistently; trend interpreted conservatively & Supporting readiness interpretation for validity and safety \\
\textbf{MET-2B-040} & Soreness & \textbf{RECOVERY} & \textbf{READINESS} & 0–10 & daily & No & \textbf{C15} & Same-day capture; scale anchored consistently; trend interpreted conservatively & Supporting readiness interpretation for validity and safety \\
\textbf{MET-2B-041} & Fatigue & \textbf{RECOVERY} & \textbf{READINESS} & 0–10 & daily & No & \textbf{C15} & Same-day capture; scale anchored consistently; trend interpreted conservatively & Supporting readiness interpretation for validity and safety \\
\textbf{MET-2B-042} & Session RPE & \textbf{RECOVERY} & \textbf{READINESS} & 0–10 & per session & No & \textbf{C15} & Same-day capture; anchored scale used; interpreted with context & Supporting session intensity governance and drift detection \\
\textbf{MET-2B-043} & Step Count & \textbf{RECOVERY} & \textbf{COMPLIANCE} & count & daily & No & \textbf{C15} & Method recorded (device/app); same-day capture; steps are baseline activity and not a session & Supporting baseline activity compliance and feasibility monitoring \\
\textbf{MET-2B-044} & Session Validity Tag & \textbf{RECOVERY} & \textbf{COMPLIANCE} & \textbf{VALID} / \textbf{MODIFIED} / \textbf{INVALID} & per session & Yes & N/A & Must be assigned same day; based on required session elements and minimum dataset; no memory-based backfill; assigned and validated via Stage 2.E conventions & Gate eligibility and admissibility computation \\
\textbf{MET-2B-045} & Cardio Minimum Standard Met & \textbf{RECOVERY} & \textbf{COMPLIANCE} & Y/N & per session & Yes & N/A & Requires modality, minutes, intensity anchor (RPE + talk-test cue), and continuity statement; derived from session record per Stage 2.E & Gate eligibility and enforcement of Operating Constraints \\
\textbf{MET-2B-046} & Weekly Admissible Session Compliance Percent & \textbf{RECOVERY} & \textbf{COMPLIANCE} & \% & weekly & Yes & N/A & Computed from admissible sessions: \textbf{VALID} plus admissible \textbf{MODIFIED} per governance; \textbf{INVALID} excluded; derived per Stage 2.E weekly audit conventions & Gate eligibility threshold input (decision window dependent) \\
\end{longtblr}

\endgroup

\newpage
\end{landscape}
\clearpage

\subsubsection*{Metric \textbf{ID}}
\phantomsection

\begin{itemize}
\item Stable identifier in \textbf{MET-2B}-\#\#\# format.
\item Used as the primary join key for thresholds, protocol cards, phase references, and crosswalks.
\end{itemize}

\subsubsection*{Metric Name}
\phantomsection

\begin{itemize}
\item Plain-language construct name.
\item Must remain stable; clarifications allowed only under controlled change.
\end{itemize}

\subsubsection*{Domain Tag}
\phantomsection

\begin{itemize}
\item Primary Domain Tag only (see 2.B.1.2).
\item Must be one of: \textbf{RUN}, \textbf{RUCK}, \textbf{SWIM}, \textbf{CAL}, \textbf{STR}, \textbf{BODYCOMP}, \textbf{DURABILITY}, \textbf{RECOVERY}, \textbf{AER}.
\end{itemize}

\subsubsection*{Measurement Type}
\phantomsection

\begin{itemize}
\item Must be one of: \textbf{TEST}, \textbf{TREND}, \textbf{READINESS}, \textbf{COMPLIANCE}.
\item Determines expected logging posture and admissibility handling.
\end{itemize}

\subsubsection*{Unit}
\phantomsection

\begin{itemize}
\item Unit-locked; changing requires controlled change.
\item Must be consistent with threshold tables and protocol cards.
\end{itemize}

\subsubsection*{Collection Cadence High-Level}
\phantomsection

\begin{itemize}
\item High-level cadence only (daily / weekly / monthly / per session / test weeks).
\item Calendar placement is owned by Stage 3; Stage 2.B does not place tests in time.
\end{itemize}

\subsubsection*{Gate-Critical}
\phantomsection

\begin{itemize}
\item Yes/No designation controlling default gate reliance posture.
\item Gate-Critical does not imply frequent testing; it implies gate relevance when tested or when computed.
\end{itemize}

\subsubsection*{Associated Protocol Card C\#}
\phantomsection

\begin{itemize}
\item Primary protocol reference that governs validity conditions and execution posture when the metric is executed as a protocolized test.
\item \textbf{TBD} 2.C permitted only in draft posture for protocolized tests.
\item N/A is required when the metric is not executed as a protocolized test (for example, derived \textbf{COMPLIANCE} metrics computed from logs under Stage 2.E).
\end{itemize}

\subsubsection*{Validity Requirements High-Level}
\phantomsection

\begin{itemize}
\item Minimum admissibility posture; must specify what must be recorded and what breaks validity.
\item Must remain consistent with Stage 1.F test validity doctrine and Stage 2.E validity condition schema.
\end{itemize}

\subsubsection*{Primary Use}
\phantomsection

\begin{itemize}
\item Why the metric exists in governance: end-state evidence, tier classification, risk control, compliance enforcement, feasibility monitoring.
\end{itemize}

Stage 2.B validity requirements must remain compatible with the Stage 2 validity condition schema tokens (\textbf{ENV}, \textbf{EQUIP}, \textbf{WARMUP}, \textbf{SAFETY}, \textbf{MEASURE}). High-level alignment rules:

\begin{itemize}
\item \textbf{TEST} metrics:
  \begin{itemize}
  \item typically require \textbf{ENV}, \textbf{EQUIP}, \textbf{WARMUP}, \textbf{SAFETY}, \textbf{MEASURE} unless the protocol explicitly permits token N/A.
  \end{itemize}
\item \textbf{TREND} metrics:
  \begin{itemize}
  \item typically require \textbf{MEASURE} and \textbf{EQUIP}; \textbf{ENV} required to establish conditions; \textbf{WARMUP} usually N/A; \textbf{SAFETY} recorded as “no stop-rule concerns” where applicable.
  \end{itemize}
\item \textbf{READINESS} metrics:
  \begin{itemize}
  \item \textbf{ENV} usually minimal; \textbf{EQUIP} N/A unless tool-based; \textbf{MEASURE} is typically “self-report method”; \textbf{WARMUP} N/A; \textbf{SAFETY} used only when the marker indicates a stop-rule concern.
  \end{itemize}
\item \textbf{COMPLIANCE} metrics:
  \begin{itemize}
  \item must be computable from logged evidence; missing fields trigger conservative handling (e.g., Cardio minimum met cannot be “assumed”).
  \end{itemize}
\end{itemize}

\subsubsection*{N/A usage posture (for validity condition tokens):}
\phantomsection

\begin{itemize}
\item N/A is permitted only when the token is truly non-applicable to the construct (for example, \textbf{WARMUP} for bodyweight measurement).
\item N/A is not permitted as a substitute for missing data on applicable tokens (for example, \textbf{MEASURE} is not N/A for run/ruck/swim distance verification).
\end{itemize}

\begin{itemize}
\item Direct metric: measured directly in a test or measurement event (e.g., run time, plank seconds, bodyweight).
\item Derived metric: computed from other records (e.g., weekly bodyweight average; weekly admissible compliance percent).
\end{itemize}

Derived metrics must define, in validity requirements, the minimum inputs required to compute them and the rule for handling missing inputs (default conservative).

\newpage

\subsection*{2.B.3 Gate-Critical Metric Set Summary}
\phantomsection
\addcontentsline{toc}{subsection}{2.B.3 Gate-Critical Metric Set Summary}

\begingroup
\scriptsize
\setlength{\tabcolsep}{2pt}
\renewcommand{\arraystretch}{1.12}

\begin{longtblr}{
  width=\linewidth,
  rowhead=1,
  colspec = {
    Q[l,wd=0.90in]
    Q[l,wd=1.00in]
    X[l,1.55]
    X[l,2.35]
    Q[l,wd=1.25in]
  },
  hlines,
  vlines,
  cells = {hsep=6pt, vsep=6pt, halign=l, valign=t},
  cell{1-Z}{1-5} = {fg=black},
  row{odd} = {bg=white},
  row{even} = {bg=LightGray},
  row{1} = {bg=StageBlue!15, fg=black, font=\bfseries, halign=l, valign=t},
}
\Hdr{Domain Tag} &
\Hdr{Metric \textbf{ID}} &
\Hdr{Metric Name} &
\Hdr{Used For Gate Context} &
\Hdr{Typical Test Window} \\
\textbf{BODYCOMP} & \textbf{MET-2B-001} & Bodyweight — Weekly Average & End-state target evidence; gate readiness posture via trend convergence & weekly \\
\textbf{BODYCOMP} & \textbf{MET-2B-002} & Body Fat Estimate Percent — Method/Device Recorded & End-state target evidence; method-consistent trend validity for readiness posture & monthly \\
\textbf{CAL} & \textbf{MET-2B-010} & Push-Ups — 2-Minute Timed, Strict & End-state target evidence; readiness tier classification input & test weeks \\
\textbf{CAL} & \textbf{MET-2B-011} & Sit-Ups — 2-Minute Timed, Standard & End-state target evidence; readiness tier classification input & test weeks \\
\textbf{CAL} & \textbf{MET-2B-012} & Pull-Ups — Strict Dead-Hang & End-state target evidence; readiness tier classification input & test weeks \\
\textbf{CAL} & \textbf{MET-2B-013} & Plank — Forearm & End-state target evidence; readiness tier classification input & test weeks \\
\textbf{RUN} & \textbf{MET-2B-014} & 1.5-Mile Run Time Trial & End-state target evidence; readiness tier classification input & test weeks \\
\textbf{RUN} & \textbf{MET-2B-015} & 2-Mile Run Time Trial & End-state target evidence; readiness tier classification input & test weeks \\
\textbf{RUN} & \textbf{MET-2B-016} & 3-Mile Run Time Trial & End-state target evidence; readiness tier classification input & test weeks \\
\textbf{RUN} & \textbf{MET-2B-017} & 5-Mile Run Time Trial & End-state target evidence; readiness tier classification input & test weeks \\
\textbf{RUCK} & \textbf{MET-2B-018} & 3-Mile Ruck Time Trial — 25 lb & End-state target evidence; readiness tier classification input & test weeks \\
\textbf{RUCK} & \textbf{MET-2B-019} & 6-Mile Ruck Time Trial — 35 lb & End-state target evidence; readiness tier classification input & test weeks \\
\textbf{RUCK} & \textbf{MET-2B-020} & 9-Mile Ruck Time Trial — 45 lb & End-state target evidence; readiness tier classification input & test weeks \\
\textbf{RUCK} & \textbf{MET-2B-021} & 12-Mile Ruck Time Trial — 55 lb & End-state target evidence; readiness tier classification input & test weeks \\
\textbf{SWIM} & \textbf{MET-2B-022} & Swim 500 yd Time — No Fins & End-state target evidence; unit-locked comparability & test weeks \\
\textbf{SWIM} & \textbf{MET-2B-023} & Swim 500 m Time — No Fins & End-state target evidence; unit-locked comparability & test weeks \\
\textbf{SWIM} & \textbf{MET-2B-025} & Underwater Exposure Admissibility Flag & Gate admissibility prerequisite for any underwater claim; governance enforcement & per underwater test event \\
\textbf{SWIM} & \textbf{MET-2B-026} & Underwater 25 m — One-Breath Time & End-state target evidence (admissible only under governance) & test weeks \\
\textbf{DURABILITY} & \textbf{MET-2B-036} & Gait-Altering Pain Flag & Gate disqualifier signal; regression/hold constraint & per session \\
\textbf{RECOVERY} & \textbf{MET-2B-044} & Session Validity Tag & Gate eligibility; admissibility computation & per session \\
\textbf{RECOVERY} & \textbf{MET-2B-045} & Cardio Minimum Standard Met & Gate eligibility; Operating Constraints enforcement (30-minute continuous cardio) & per session \\
\textbf{RECOVERY} & \textbf{MET-2B-046} & Weekly Admissible Session Compliance Percent & Gate eligibility threshold input; decision window compliance posture & weekly \\
\end{longtblr}

\endgroup

\begin{itemize}
\item If a Gate-Critical metric is missing for a gate window where it is required:
  \begin{itemize}
  \item default outcome is \textbf{HOLD} unless safety posture requires \textbf{REGRESS} or \textbf{MEDICAL} \textbf{HOLD}.
  \end{itemize}
\item If a Gate-Critical metric is present but validity conditions are not satisfied:
  \begin{itemize}
  \item evidence is treated as non-admissible for gate purposes,
  \item corrective action is re-execution under valid conditions (for tests) or restoration of method consistency (for trends).
  \end{itemize}
\item Gate-Critical does not mean “test more frequently.” It means “when the program claims readiness or progression, this evidence must exist and be admissible.”
\end{itemize}

\newpage

\subsection*{2.B.4 Metric-to-Protocol Coverage Map}
\phantomsection
\addcontentsline{toc}{subsection}{2.B.4 Metric-to-Protocol Coverage Map}

\begingroup
\scriptsize
\setlength{\tabcolsep}{2pt}
\renewcommand{\arraystretch}{1.12}

\begin{longtblr}{
  width=\linewidth,
  rowhead=1,
  colspec = {
    Q[l,wd=1.00in]
    X[l,2.45]
    Q[l,wd=1.35in]
    Q[l,wd=0.95in]
  },
  hlines,
  vlines,
  cells = {hsep=6pt, vsep=6pt, halign=l, valign=t},
  cell{1-Z}{1-4} = {fg=black},
  row{odd} = {bg=white},
  row{even} = {bg=LightGray},
  row{1} = {bg=StageBlue!15, fg=black, font=\bfseries, halign=l, valign=t},
}
\Hdr{Metric \textbf{ID}} &
\Hdr{Metric Name} &
\Hdr{Intended Protocol Card C\#} &
\Hdr{Status} \\
\textbf{MET-2B-001} & Bodyweight — Weekly Average & \textbf{C1} & Accepted/Complete \\
\textbf{MET-2B-002} & Body Fat Estimate Percent — Method/Device Recorded & \textbf{C2} & Accepted/Complete \\
\textbf{MET-2B-003} & Circumference — Neck & \textbf{C3} & Accepted/Complete \\
\textbf{MET-2B-004} & Circumference — Chest & \textbf{C3} & Accepted/Complete \\
\textbf{MET-2B-005} & Circumference — Bicep & \textbf{C3} & Accepted/Complete \\
\textbf{MET-2B-006} & Circumference — Forearm & \textbf{C3} & Accepted/Complete \\
\textbf{MET-2B-007} & Circumference — Waist & \textbf{C3} & Accepted/Complete \\
\textbf{MET-2B-008} & Circumference — Thigh & \textbf{C3} & Accepted/Complete \\
\textbf{MET-2B-009} & Circumference — Calf & \textbf{C3} & Accepted/Complete \\
\textbf{MET-2B-010} & Push-Ups — 2-Minute Timed, Strict & \textbf{C4} & Accepted/Complete \\
\textbf{MET-2B-011} & Sit-Ups — 2-Minute Timed, Standard & \textbf{C5} & Accepted/Complete \\
\textbf{MET-2B-012} & Pull-Ups — Strict Dead-Hang & \textbf{C6} & Accepted/Complete \\
\textbf{MET-2B-013} & Plank — Forearm & \textbf{C7} & Accepted/Complete \\
\textbf{MET-2B-014} & 1.5-Mile Run Time Trial & \textbf{C8} & Accepted/Complete \\
\textbf{MET-2B-015} & 2-Mile Run Time Trial & \textbf{C9} & Accepted/Complete \\
\textbf{MET-2B-016} & 3-Mile Run Time Trial & \textbf{C10} & Accepted/Complete \\
\textbf{MET-2B-017} & 5-Mile Run Time Trial & \textbf{C11} & Accepted/Complete \\
\textbf{MET-2B-018} & 3-Mile Ruck Time Trial — 25 lb & \textbf{C12} & Accepted/Complete \\
\textbf{MET-2B-019} & 6-Mile Ruck Time Trial — 35 lb & \textbf{C12} & Accepted/Complete \\
\textbf{MET-2B-020} & 9-Mile Ruck Time Trial — 45 lb & \textbf{C12} & Accepted/Complete \\
\textbf{MET-2B-021} & 12-Mile Ruck Time Trial — 55 lb & \textbf{C12} & Accepted/Complete \\
\textbf{MET-2B-022} & Swim 500 yd Time — No Fins & \textbf{C13} & Accepted/Complete \\
\textbf{MET-2B-023} & Swim 500 m Time — No Fins & \textbf{C13} & Accepted/Complete \\
\textbf{MET-2B-024} & Surface Water Confidence Exposure Flag & \textbf{C13} & Accepted/Complete \\
\textbf{MET-2B-025} & Underwater Exposure Admissibility Flag & \textbf{C14} & Accepted/Complete \\
\textbf{MET-2B-026} & Underwater 25 m — One-Breath Time & \textbf{C14} & Accepted/Complete \\
\textbf{MET-2B-027} & Hinge Strength Proxy — \textbf{3-RM} / \textbf{5-RM} Load & \textbf{C17} & Accepted/Complete \\
\textbf{MET-2B-028} & Squat Strength Proxy — \textbf{3-RM} / \textbf{5-RM} Load & \textbf{C17} & Accepted/Complete \\
\textbf{MET-2B-029} & Upper Push Strength Proxy — \textbf{3-RM} / \textbf{5-RM} Load & \textbf{C17} & Accepted/Complete \\
\textbf{MET-2B-030} & Upper Pull Strength Proxy — \textbf{3-RM} / \textbf{5-RM} Load & \textbf{C17} & Accepted/Complete \\
\textbf{MET-2B-031} & Carry Capacity — Load Outcome & \textbf{C17} & Accepted/Complete \\
\textbf{MET-2B-032} & Rower Time Trial & \textbf{C18} & Accepted/Complete \\
\textbf{MET-2B-033} & Bike Time Trial & \textbf{C18} & Accepted/Complete \\
\textbf{MET-2B-034} & Elliptical Time Trial & \textbf{C18} & Accepted/Complete \\
\textbf{MET-2B-035} & Pain Severity — 0–10 + Location Recorded & \textbf{C16} & Accepted/Complete \\
\textbf{MET-2B-036} & Gait-Altering Pain Flag & \textbf{C16} & Accepted/Complete \\
\textbf{MET-2B-037} & Foot Hotspot or Blister Flag + Location & \textbf{C16} & Accepted/Complete \\
\textbf{MET-2B-038} & Sleep Hours & \textbf{C15} & Accepted/Complete \\
\textbf{MET-2B-039} & Sleep Quality & \textbf{C15} & Accepted/Complete \\
\textbf{MET-2B-040} & Soreness & \textbf{C15} & Accepted/Complete \\
\textbf{MET-2B-041} & Fatigue & \textbf{C15} & Accepted/Complete \\
\textbf{MET-2B-042} & Session RPE & \textbf{C15} & Accepted/Complete \\
\textbf{MET-2B-043} & Step Count & \textbf{C15} & Accepted/Complete \\
\textbf{MET-2B-044} & Session Validity Tag & N/A & Accepted/Complete \\
\textbf{MET-2B-045} & Cardio Minimum Standard Met & N/A & Accepted/Complete \\
\textbf{MET-2B-046} & Weekly Admissible Session Compliance Percent & N/A & Accepted/Complete \\
\end{longtblr}

\endgroup

\begin{itemize}
\item “Drafted” means the protocol card content exists in Stage 2.C for protocolized tests, or the metric is governed as a derived/log-computed construct under Stage 2.E when Intended Protocol Card C\# is N/A.
\item “Draft planned” means the protocol card \textbf{ID} is assigned and expected to govern the metric, but the protocol card content may still be authored in Stage 2.C.
\item “Not drafted” means protocol assignment and/or protocol content remains incomplete; metrics with “Not drafted” protocol assignment must not be treated as gate-grade unless they are non-gate-critical supporting markers used informationally.
\item Gate-Critical metrics listed as \textbf{TBD} 2.C are permitted only while Stage 2 is in Draft posture. For Stage 2 to be Deployable, Gate-Critical metrics must have protocol card coverage assigned and authored, or be explicitly governed as N/A (derived/log-computed) with Stage 2.E defining the computation and admissibility posture.
\end{itemize}

\newpage

\subsection*{2.B.5 Conflict-Resolution Rule}
\phantomsection
\addcontentsline{toc}{subsection}{2.B.5 Conflict-Resolution Rule}

If any downstream artifact references an undefined \textbf{MET-2B}-\#\#\#, uses non-canonical domain tags, or treats inadmissible data as gate evidence, the downstream artifact is Draft — Non-Deployable until corrected and re-versioned.

A downstream artifact \textbf{MUST} be labeled Draft — Non-Deployable if any of the following occurs:

\begin{itemize}
\item References a \textbf{MET-2B}-\#\#\# \textbf{ID} not present in the Stage 2.B inventory.
\item Assigns more than one Primary Domain Tag to a single metric/test construct.
\item Treats \textbf{AER} evidence as \textbf{RUN} evidence or uses \textbf{AER} constructs to claim \textbf{RUN} tolerance/performance.
\item Uses invalid or unverified measurement conditions for gate-critical tests (unknown distance, unknown load, unknown swim unit) and still claims the result is admissible for gates.
\item Attempts to use compliance metrics (Session Validity Tag, Cardio Minimum Standard Met, Weekly Admissible Session Compliance Percent) without the underlying minimum dataset required to compute them.
\end{itemize}

\begin{itemize}
\item Correction must explicitly identify:
  \begin{itemize}
  \item the offending reference or tag,
  \item the corrected \textbf{ID}/tag,
  \item the date of correction,
  \item the owner responsible.
  \end{itemize}
\item If the error affected gate decisions, an evidence integrity review is required:
  \begin{itemize}
  \item determine whether the gate decision relied on the incorrect metric/tag,
  \item if yes, the gate decision must be treated as non-authoritative until rerun using admissible evidence under correct IDs and tags.
  \end{itemize}
\end{itemize}

\end{document}